\documentclass{article}
\usepackage[margin=2cm]{geometry}
\usepackage{amsmath}
\usepackage{graphicx}
\usepackage{booktabs}
\usepackage[utf8]{inputenc}
\usepackage{ragged2e}
\setlength{\emergencystretch}{3em}
\begin{document}
\sloppy

\noindent\hfill\includegraphics[height=1.8em]{assets/figure_header_logo.png}\par\noindent\hrule\vspace{4pt}

\subsection*{1 Proof of Theorem 3}


\subsection*{1.1 Average direct effect on the never-takers}


In the following, we show that 0 = E[Yi(1, 0)-Y(0, 0)] = n] = E[Y-Qoo(Yo)|D =
1, M = 0]. From O), we obtain the first ingredient E[Yi(1,0)] = n] = E[Yi|D =
1, M = 0]. Furthermore, from () we have E[Qoo(Yo)|D = 1, M = 0] = E[Yi(0,0)|D =
1, M(1) = 0]. Under Assumption 7 and 8,


\begin{equation}
\begin{array} { r } { E [ Y _ { 1 } ( 0, 0 ) | D = 1, M ( 1 ) = 0 ] \stackrel { A 7 } { = } E [ Y _ { 1 } ( 0, 0 ) | D = 1, \tau = n ] \qquad \qquad \qquad ( 1 ) } \\ { \stackrel { \Delta 8 } { = } E [ Y _ { 1 } ( 0, 0 ) | \tau = n ]. \qquad \qquad ( 1 ) } \end{array}
\end{equation}


1.2  Quantile direct effect on the never-takers
We prove that


\begin{equation}
\begin{array} { r } { \theta _ { 1 } ^ { n } ( q ) = F _ { Y _ { 1 } ( 1, 0 ) | \tau = n } ^ { - 1 } ( q ) - F _ { Y _ { 1 } ( 0, 0 ) | \tau = n } ^ { - 1 } ( q ), } \\ { = F _ { Y _ { 1 } | D = 1, M = 0 } ^ { - 1 } ( q ) - F _ { Q _ { 0 0 } ( Y _ { 0 } ) | D = 1, M = 0 } ^ { - 1 } ( q ). } \end{array}
\end{equation}


This requires showing that


\begin{equation}
\begin{array} { c } { F _ { Y _ { 1 } ( 1, 0 ) | \tau = n } ( y ) = F _ { Y _ { 1 } | D = 1, M = 0 } ( y ) ~ \mathrm { a n d } \qquad \qquad ( 2 ) } \\ { F _ { Y _ { 1 } ( 0, 0 ) | \tau = n } ( y ) = F _ { Q _ { 0 0 } ( Y _ { 0 } ) | D = 1, M = 0 } ( y ). \qquad \qquad \qquad ( 3 ) } \end{array}
\end{equation}


Under Assumptions 7 and 8.


\begin{equation}
\begin{array} { c } { F _ { Y _ { t } | D = 1, M = 0 } ( y ) = E [ 1 \{ Y _ { t } \leq y \} | D = 1, M = 0 ] } \\ { \overset { A. A 8 } { = } E [ 1 \{ Y _ { t } ( 1, 0 ) \leq y \} | \tau = n ] \qquad \qquad ( 4 ) } \\ { = F _ { Y _ { t } ( 1, 0 ) | \tau = n } ( y ), \qquad \qquad \qquad \qquad ( 4 ) } \end{array}
\end{equation}


which proves O. From O, we have


FQoo(Yo)ID=1,M=0(y) = FY(0.0)ID=1,M(1)=0(y) = E[1\{Yi(0,0)   y\}|D = 1, M(1) = 0].

\end{document}
